\section{Additional Comparison Results}
\label{sec:results}
\subsection{Task-to-PCC response}
\label{sec:freq}
The band error of a task depends on the halls that host it, as \eqref{eq:task-error} shows.
Table~\ref{tab:task-band} therefore gives the active-power comparison for every task and both realizations.
The inference task appears because its allocation defines a task-to-PCC response.
Its power is constant within each window, so it does not excite this response in the time-domain runs.
Its resonance-band error ranges from 5.63\% to 5.83\%.
These values are close to the band error of the plant-wide job of S1, since the inference nodes also occupy every hall.
{\small\setlength{\tabcolsep}{4.5pt}\setlength{\LTcapwidth}{\textwidth}
\begin{longtable}{lllrrrrrrr}
\caption{Task-to-PCC Active-Power Comparison for Every Task and Both Realizations.
Paper Denotes the Realization Used in the Paper.
The Six Jobs of S2 Are Training Jobs.}\label{tab:task-band}\\
\toprule
 & & & \multicolumn{2}{c}{Band error, \%} & \multicolumn{2}{c}{Full model peak} & \multicolumn{2}{c}{Equivalent peak} & Phase\\
Scenario & Task & Realization & 0.1--3 Hz & 3--20 Hz & Gain & Hz & Gain & Hz & deg\\\midrule
\endfirsthead
\toprule
 & & & \multicolumn{2}{c}{Band error, \%} & \multicolumn{2}{c}{Full model peak} & \multicolumn{2}{c}{Equivalent peak} & Phase\\
Scenario & Task & Realization & 0.1--3 Hz & 3--20 Hz & Gain & Hz & Gain & Hz & deg\\\midrule
\endhead
S1 & training & paper & 1.33 & 5.82 & 1.904 & 5.33 & 1.978 & 5.49 & 3.2 \\*
 &  & second & 1.32 & 5.74 & 1.891 & 5.31 & 1.964 & 5.47 & 3.1 \\
S1 & inference & paper & 1.33 & 5.83 & 1.904 & 5.33 & 1.978 & 5.49 & 3.2 \\*
 &  & second & 1.32 & 5.74 & 1.891 & 5.31 & 1.964 & 5.47 & 3.1 \\
\midrule
S2 & job 1 & paper & 5.18 & 27.84 & 1.872 & 4.83 & 1.978 & 5.49 & 15.2 \\*
 &  & second & 5.14 & 27.45 & 1.859 & 4.80 & 1.963 & 5.47 & 14.9 \\
S2 & job 2 & paper & 3.10 & 21.33 & 2.143 & 6.24 & 1.978 & 5.49 & $-$13.9 \\*
 &  & second & 3.06 & 21.02 & 2.120 & 6.20 & 1.963 & 5.47 & $-$13.7 \\
S2 & job 3 & paper & 0.95 & 13.86 & 2.174 & 6.02 & 1.978 & 5.49 & 0.1 \\*
 &  & second & 0.92 & 13.33 & 2.147 & 5.97 & 1.963 & 5.47 & $-$0.1 \\
S2 & job 4 & paper & 4.56 & 24.54 & 1.901 & 5.01 & 1.978 & 5.49 & 13.4 \\*
 &  & second & 4.53 & 24.19 & 1.886 & 4.99 & 1.963 & 5.47 & 13.2 \\
S2 & job 5 & paper & 0.76 & 10.57 & 1.876 & 5.09 & 1.978 & 5.49 & 1.0 \\*
 &  & second & 0.75 & 10.27 & 1.867 & 5.08 & 1.963 & 5.47 & 1.0 \\
S2 & job 6 & paper & 0.54 & 7.50 & 1.871 & 5.50 & 1.978 & 5.49 & $-$2.8 \\*
 &  & second & 0.54 & 7.30 & 1.860 & 5.48 & 1.963 & 5.47 & $-$2.7 \\
S2 & inference & paper & 1.33 & 5.83 & 1.904 & 5.33 & 1.978 & 5.49 & 3.2 \\*
 &  & second & 1.32 & 5.73 & 1.891 & 5.31 & 1.963 & 5.47 & 3.1 \\
\midrule
S3 & training & paper & 1.82 & 9.55 & 1.893 & 5.19 & 1.979 & 5.50 & 5.5 \\*
 &  & second & 1.81 & 9.41 & 1.879 & 5.17 & 1.963 & 5.47 & 5.4 \\
S3 & fine-tuning & paper & 0.98 & 4.74 & 1.918 & 5.33 & 1.979 & 5.50 & 2.7 \\*
 &  & second & 0.97 & 4.65 & 1.904 & 5.31 & 1.963 & 5.47 & 2.6 \\
S3 & inference & paper & 1.33 & 5.72 & 1.905 & 5.34 & 1.979 & 5.50 & 3.1 \\*
 &  & second & 1.32 & 5.63 & 1.891 & 5.31 & 1.963 & 5.47 & 3.0 \\
\bottomrule\end{longtable}}


Table~\ref{tab:output-band} gives the largest band error over the jobs for each of the four outputs.
In the realization used in the paper, the largest value in the resonance band is 7.38\%, 30.47\% and 11.57\% for S1, S2 and S3.
The reactive power gives this value in S1 and S3, and the PCC voltage in S2.
The second realization changes the band error of every task, output and band by at most 0.61 percentage points.
\begin{table}[H]
\centering\small\setlength{\tabcolsep}{4.5pt}
\caption{Largest Band Error over the Jobs of Each Realization for the Four Outputs, in Percent.
Bands in Hz.}
\label{tab:output-band}
\begin{tabular}{llrrrrrrrr}\toprule
 & & \multicolumn{2}{c}{Active power} & \multicolumn{2}{c}{Reactive power} & \multicolumn{2}{c}{PCC voltage} & \multicolumn{2}{c}{Machine frequency}\\
Scenario & Realization & 0.1--3 & 3--20 & 0.1--3 & 3--20 & 0.1--3 & 3--20 & 0.1--3 & 3--20\\\midrule
S1 & paper & 1.33 & 5.82 & 1.58 & 7.38 & 1.09 & 6.35 & 0.90 & 5.10 \\
S1 & second & 1.32 & 5.74 & 1.57 & 7.28 & 1.08 & 6.24 & 0.89 & 5.02 \\
S2 & paper & 5.18 & 27.84 & 13.22 & 27.52 & 6.09 & 30.47 & 3.32 & 25.34 \\
S2 & second & 5.14 & 27.45 & 13.18 & 27.11 & 6.01 & 30.02 & 3.29 & 24.99 \\
S3 & paper & 1.82 & 9.55 & 2.48 & 11.57 & 1.57 & 10.61 & 1.17 & 8.44 \\
S3 & second & 1.81 & 9.41 & 2.48 & 11.41 & 1.56 & 10.44 & 1.16 & 8.32 \\
\bottomrule\end{tabular}\end{table}


Table~\ref{tab:hall-band} gives the single-hall and uniform-input band errors.
Over the six inputs, a unit input in one hall gives band errors from 5.46\% to 53.17\%, and the uniform input gives 4.57\% to 4.64\%.
Hall~1 gives the largest and hall~12 the smallest single-hall error at every operating point.
These errors depend on the scenario only through its operating point, so the six columns nearly coincide.
\begin{table}[H]
\centering\small
\caption{PCC Active-Power Band Error over 3--20 Hz, in Percent, for a Unit Input in One Hall and for the Same Input Spread Evenly over All Halls}
\label{tab:hall-band}
\begin{tabular}{lrrrrrr}\toprule
 & \multicolumn{2}{c}{S1} & \multicolumn{2}{c}{S2} & \multicolumn{2}{c}{S3}\\
Input & Paper & Second & Paper & Second & Paper & Second\\\midrule
1 & 53.17 & 52.66 & 53.16 & 52.66 & 53.16 & 52.64 \\
2 & 6.61 & 6.59 & 6.61 & 6.59 & 6.67 & 6.65 \\
3 & 19.26 & 18.88 & 19.24 & 18.88 & 19.34 & 18.92 \\
4 & 26.28 & 25.92 & 26.27 & 25.91 & 26.31 & 25.92 \\
5 & 9.19 & 8.94 & 9.19 & 8.89 & 9.26 & 8.99 \\
6 & 19.69 & 19.07 & 19.69 & 18.99 & 19.75 & 19.05 \\
7 & 38.83 & 38.40 & 38.83 & 38.38 & 38.76 & 38.30 \\
8 & 16.10 & 15.78 & 16.10 & 15.77 & 16.16 & 15.80 \\
9 & 11.24 & 11.02 & 11.24 & 11.00 & 11.33 & 11.06 \\
10 & 14.05 & 13.69 & 14.04 & 13.67 & 14.05 & 13.66 \\
11 & 18.78 & 18.23 & 18.76 & 18.21 & 18.79 & 18.21 \\
12 & 5.59 & 5.47 & 5.60 & 5.46 & 5.55 & 5.47 \\
\midrule
uniform & 4.64 & 4.57 & 4.64 & 4.57 & 4.64 & 4.57 \\
\bottomrule\end{tabular}\end{table}


Table~\ref{tab:modes} lists the modes between 3 and 20~Hz that reach 10\% of the largest contribution of the full model.
The full model has 13 to 15 such modes, whereas the equivalent has one in S2 and three in S1 and S3.
\begin{table}[H]
\centering\small\setlength{\tabcolsep}{4.5pt}
\caption{Modes between 3 and 20 Hz That Reach 10\% of the Largest Contribution of the Full Model}
\label{tab:modes}
\begin{tabular}{llrllrll}\toprule
 & & \multicolumn{3}{c}{Full model} & \multicolumn{3}{c}{Single-block equivalent}\\
Scenario & Realization & Modes & Frequency, Hz & Damping ratio & Modes & Frequency, Hz & Damping ratio\\\midrule
S1 & paper & 15 & 5.01--14.52 & 0.336--0.727 & 3 & 5.12, 6.03, 14.98 & 0.336, 0.458, 0.757 \\
S1 & second & 15 & 4.99--14.49 & 0.337--0.726 & 3 & 5.11, 5.99, 15.02 & 0.337, 0.462, 0.756 \\
S2 & paper & 13 & 5.01--7.00 & 0.336--0.520 & 1 & 6.03 & 0.458 \\
S2 & second & 13 & 4.99--6.96 & 0.337--0.522 & 1 & 5.99 & 0.462 \\
S3 & paper & 14 & 5.01--14.31 & 0.336--0.727 & 3 & 5.12, 6.03, 14.97 & 0.336, 0.457, 0.757 \\
S3 & second & 14 & 4.99--14.37 & 0.337--0.726 & 3 & 5.11, 5.99, 15.02 & 0.337, 0.462, 0.756 \\
\bottomrule\end{tabular}\end{table}


\subsection{Time-domain response}
\label{sec:time}
Table~\ref{tab:dynamic} gives the dynamic errors of the four outputs for both realizations, the initial offsets and the spectral share of the active-power error.
The dynamic error of the PCC active power ranges from 1.46\% to 5.94\%, and S2 gives the largest error of every output in both realizations.
The equivalent starts with a PCC active power between 6.41 and 7.72~kW below that of the full model, because its losses at the operating point differ from those of the full model.
The initial frequency offset is zero in every run.
Between 73.7\% and 82.7\% of the variance of the active-power error lies between 3 and 9~Hz, inside the resonance band.
\begin{table}[H]
\centering\small\setlength{\tabcolsep}{4.5pt}
\caption{Dynamic Errors of the Four Outputs, Initial Offsets and Share of the Active-Power Error Variance between 3 and 9 Hz}
\label{tab:dynamic}
\begin{tabular}{llrrrrrrrr}\toprule
 & & \multicolumn{4}{c}{Dynamic error, \%} & \multicolumn{3}{c}{Initial offset $b_0$} & Share\\
Scenario & Realization & $P$ & $Q$ & $V$ & $f$ & $P$, kW & $Q$, kvar & $V$, $10^{-6}$ p.u. & 3--9 Hz, \%\\\midrule
S1 & paper & 1.62 & 2.08 & 1.80 & 0.46 & $-$6.93 & $-$0.60 & 0.99 & 73.7 \\
S1 & second & 2.41 & 3.10 & 3.30 & 0.74 & $-$7.18 & $-$1.29 & 1.52 & 79.2 \\
S2 & paper & 5.69 & 8.92 & 7.16 & 1.25 & $-$6.41 & 4.48 & $-$2.70 & 82.3 \\
S2 & second & 5.94 & 8.21 & 7.54 & 1.38 & $-$7.72 & $-$2.13 & 2.17 & 82.6 \\
S3 & paper & 2.10 & 2.77 & 2.98 & 0.57 & $-$7.66 & $-$0.92 & 1.30 & 82.7 \\
S3 & second & 1.46 & 1.94 & 1.73 & 0.34 & $-$7.23 & $-$2.58 & 2.44 & 79.2 \\
\bottomrule\end{tabular}\end{table}


A strict rerun checks whether the production tolerance affects these errors.
It repeats both models for S2 in the realization used in the paper at relative tolerance $10^{-8}$ and absolute tolerance $10^{-10}$.
The dynamic error of the PCC active power is then 5.6929\%, against 5.6924\% in the production run, a difference of 0.0005 percentage points.
The largest change of the other three outputs is 0.0018 percentage points.

\subsection{Internal variables}
\label{sec:internal}
The paper shows the internal variables of S2 in its Fig.~4.
Table~\ref{tab:internal} gives their peak deviations for all three scenarios in the realization used in the paper.
For every block variable, the equivalent deviates less than the most strongly excited hall of the full model.
The ratio ranges from 0.19 to 0.37 in S2 and from 0.56 to 0.81 in S1 and S3.
In S2, the equivalent also deviates less than every hall, because the six independent jobs partly cancel in the total server power that drives it.
In S1 and S3, its peak deviations lie within the range of the halls.
The largest peak deviation of the AC input power among the halls exceeds that of the server power by a factor of 1.5 to 1.7 in every scenario.
\begin{table}[tbp]
\centering\small\setlength{\tabcolsep}{4.5pt}
\caption{Peak Deviations of Internal Variables in the Realization Used in the Paper.
Per-Unit Values on the Block Base.}
\label{tab:internal}
\begin{tabular}{lllrrlrr}\toprule
 & & & \multicolumn{3}{c}{Full model} & & \\
Scenario & Quantity & Unit & Hall & Largest & Range over halls & Equivalent & Ratio\\\midrule
S1 & Server power $P_{{\rm IT},h}$ & MW & 1 & 0.364 & 0.228--0.364 & --- & --- \\
 & AC input power $P_{{\rm AC},h}$ & MW & 1 & 0.594 & 0.360--0.594 & --- & --- \\
 & DC-link voltage $v_d$ & p.u. & 1 & 0.1159 & 0.0453--0.1159 & 0.0711 & 0.613 \\
 & VSI output voltage $|\bm v_v|$ & p.u. & 6 & 0.0107 & 0.0042--0.0107 & 0.0068 & 0.633 \\
 & PSU output voltage $v_p$ & p.u. & 1 & 0.0419 & 0.0201--0.0419 & 0.0272 & 0.649 \\
 & Server voltage $v_e$ & p.u. & 8 & 0.0130 & 0.0081--0.0130 & 0.0105 & 0.807 \\
 & AFE current $|\bm i_a|$ & p.u. & 1 & 0.0518 & 0.0317--0.0518 & 0.0392 & 0.757 \\
\midrule
S2 & Server power $P_{{\rm IT},h}$ & MW & 8 & 0.341 & 0.183--0.341 & --- & --- \\
 & AC input power $P_{{\rm AC},h}$ & MW & 1 & 0.569 & 0.310--0.569 & --- & --- \\
 & DC-link voltage $v_d$ & p.u. & 1 & 0.1136 & 0.0411--0.1136 & 0.0274 & 0.241 \\
 & VSI output voltage $|\bm v_v|$ & p.u. & 6 & 0.0105 & 0.0041--0.0105 & 0.0024 & 0.228 \\
 & PSU output voltage $v_p$ & p.u. & 1 & 0.0430 & 0.0213--0.0430 & 0.0100 & 0.233 \\
 & Server voltage $v_e$ & p.u. & 1 & 0.0134 & 0.0086--0.0134 & 0.0026 & 0.190 \\
 & AFE current $|\bm i_a|$ & p.u. & 1 & 0.0487 & 0.0263--0.0487 & 0.0180 & 0.369 \\
\midrule
S3 & Server power $P_{{\rm IT},h}$ & MW & 2 & 0.448 & 0.168--0.448 & --- & --- \\
 & AC input power $P_{{\rm AC},h}$ & MW & 2 & 0.652 & 0.213--0.652 & --- & --- \\
 & DC-link voltage $v_d$ & p.u. & 1 & 0.1243 & 0.0384--0.1243 & 0.0724 & 0.583 \\
 & VSI output voltage $|\bm v_v|$ & p.u. & 7 & 0.0118 & 0.0035--0.0118 & 0.0068 & 0.578 \\
 & PSU output voltage $v_p$ & p.u. & 7 & 0.0490 & 0.0164--0.0490 & 0.0276 & 0.563 \\
 & Server voltage $v_e$ & p.u. & 2 & 0.0158 & 0.0053--0.0158 & 0.0095 & 0.603 \\
 & AFE current $|\bm i_a|$ & p.u. & 2 & 0.0560 & 0.0190--0.0560 & 0.0316 & 0.564 \\
\bottomrule\end{tabular}\end{table}


Fig.~\ref{fig:internal} shows S1 and S3 in the layout of Fig.~4 of the paper.
Each column uses the window registered for its scenario from the input alone.
The window begins 0.5~s before the phase transition with the largest step of the total server power and ends 1.5~s after it, widened to the nearest axis ticks.
This step is $-3.67$~MW at 9.42~s in S1 and $-2.58$~MW at 36.23~s in S3.
Each phase transition is a step of the server power in the halls that host the job.
The step starts a damped oscillation of their AC input power, DC-link voltage and PSU output voltage.
In both scenarios, the PCC active power of the two models nearly coincides, whereas the DC-link and PSU output voltages of the halls spread around the one trace of the equivalent.
\begin{figure}[tbp]
\centering\includegraphics[width=\textwidth]{internal_variables.pdf}
\caption{Internal variables of S1 in the left column and of S3 in the right column, in the realization used in the paper.
(a) Server power $P_{{\rm IT},h}$ and (b) AC input power $P_{{\rm AC},h}$ of every hall.
(c) Change of the PCC active power.
(d) DC-link voltage $v_d$, (e) PSU output voltage $v_p$ and (f) server voltage $v_e$ of every hall.
Black lines show the PCC or the hall with the largest DC-link excursion of the full model in the shown window, which is hall~1 in S1 and hall~2 in S3.
Grey lines show the other halls, and dashed lines the single-block equivalent.}
\label{fig:internal}
\end{figure}

\subsection{Numerical checks}
\label{sec:checks}
The single-block equivalent passes its registered gates at all six operating points.
The equivalent has 179 equations and variables, and the full model has 1873.
Over the six operating points, the largest initial residual is $4.0\cdot10^{-12}$ for the equivalent and $5.8\cdot10^{-12}$ for the full model.
The server power of the equivalent equals that of the full model and of the scenario within $5.8\cdot10^{-12}$~MW, and its static loads equal those of the full model.
A two-second run without input checks that each operating point is an equilibrium.
The largest change of any variable of the equivalent is $5.1\cdot10^{-13}$, and that of the full model is $4.5\cdot10^{-13}$ at the three operating points of the realization used in the paper.
Both models are stable at all six operating points apart from the zero mode of the common rotor angle.
The largest remaining real part of the full model is $-0.117$~s$^{-1}$.
In the positive control of Section~\ref{sec:equivalent}, each eigenvalue of the equivalent is compared with the nearest eigenvalue of the full model.
Their distance is divided by one plus the magnitude of the eigenvalue of the equivalent.
The largest normalized distance is $9.1\cdot10^{-12}$.

The internal-variable reruns reproduce the production runs of both models, with zero difference in the PCC active power and in the DC-link voltage at every output instant.
The four blocks of every hall of the full model agree in every recorded variable within $4.4\cdot10^{-14}$.
The six task-driven runs of the full model keep the converter current below 0.690~p.u., every modulation index below 1.143 times its operating-point value and the DC-link voltage between 0.875 and 1.113~p.u.

\subsection{Evidence limits}
One construction of the single-block equivalent was tested, namely the capacity-weighted average block of Section~\ref{sec:equivalent}.
A single nominal block scaled to the plant rating was not tested.
The comparisons use one electrical design of the 12 halls.
Its parameter ranges are study assumptions and do not form a measured distribution of data-center equipment.

The fast inputs combine literature timing with measured node-power levels.
They are constructed scenarios.
They do not constitute measured PCC transients, measured synchronization across an operating facility or a high-bandwidth measurement of the NLR training cycles.
The polling intervals of the NLR records do not establish the sensor bandwidth.
The phase parameters are the experimental settings of \cite{ko2026oscillations} and not universal workload limits.
Inference remains at a constant request rate within each task window.
The selected job size of 0.125 of the baseline job slots is not a maximum safe job size for a physical facility.

The averaged supply-block model excludes switching ripple, current and modulation limits, protection, battery discharge and bypass switching.
The checks of Section~\ref{sec:checks} therefore describe the response of the model to the selected inputs and do not validate converter protection or field operating limits.
The quasi-static network represents electromagnetic line behavior only through fixed-frequency impedances.
A network with differential electromagnetic states is required to test the resulting approximation across the full 20~Hz band.
The present conclusions concern task-driven PCC quantities and block variables under the stated conditions.
